% ============================================================
% oreilly-style.tex — Stile O'Reilly Media per libri tecnici
% ============================================================
% Caratteristiche: font serif elegante, capitoli con animale
% illustrativo, layout pulito e professionale, colori sobri.
% ============================================================

% === LINGUA ===
\usepackage[\BabelLang]{babel}
\usepackage[autostyle]{csquotes}

% === BIBLIOGRAFIA ===
\usepackage[
  backend=biber,
  style=numeric,
  sorting=nyt,
  urldate=long,
  maxbibnames=3,
  backref=true,
  backrefstyle=none
]{biblatex}

% === FONT (LuaLaTeX — stile O'Reilly: Garamond + Helvetica) ===
\usepackage{fontspec}
\usepackage{amsmath}
\usepackage{amssymb}
\setmainfont{texgyretermes}[
  Extension=.otf,
  UprightFont=*-regular,
  BoldFont=*-bold,
  ItalicFont=*-italic,
  BoldItalicFont=*-bolditalic
]
\setsansfont{texgyreheros}[
  Extension=.otf,
  Scale=0.95,
  UprightFont=*-regular,
  BoldFont=*-bold,
  ItalicFont=*-italic,
  BoldItalicFont=*-bolditalic
]
\setmonofont{InconsolataN}[
  Extension=.otf,
  Scale=0.85,
  UprightFont=*-Regular,
  BoldFont=*-Bold
]

% === LAYOUT PAGINA (7 × 9.1875 poll. — formato O'Reilly) ===
\usepackage[
  paperwidth=7in,
  paperheight=9.1875in,
  top=0.875in, bottom=0.75in,
  inner=1in, outer=0.75in,
  headheight=14pt,
  footskip=0.4in
]{geometry}

% === SPAZIATURA PARAGRAFI (stile O'Reilly: indent, no extra space) ===
\setlength{\parindent}{1.5em}
\setlength{\parskip}{0pt plus 1pt}

% === COLORI (palette O'Reilly — teal/grigio) ===
\usepackage[dvipsnames,svgnames,table]{xcolor}
\definecolor{oreilly-teal}{HTML}{006B6B}
\definecolor{oreilly-dark}{HTML}{1A1A1A}
\definecolor{oreilly-gray}{HTML}{555555}
\definecolor{oreilly-lightgray}{HTML}{999999}
\definecolor{code-bg}{HTML}{F7F7F7}
\definecolor{code-border}{HTML}{CCCCCC}
\definecolor{tip-green}{HTML}{528B2E}
\definecolor{warn-orange}{HTML}{CC6600}
\definecolor{note-blue}{HTML}{336699}
\definecolor{link-blue}{HTML}{004B6B}
\definecolor{caution-red}{HTML}{993333}

% === GRAFICA E TABELLE ===
\usepackage{graphicx}
\usepackage{float}
\usepackage{booktabs}
\usepackage{longtable}
\usepackage{tabularx}
\usepackage{multirow}
\usepackage{array}
\usepackage{colortbl}

% === CODICE SORGENTE (stile O'Reilly: box leggero, bordo sottile) ===
\usepackage{listings}
\lstset{
  basicstyle=\footnotesize\ttfamily\color{oreilly-dark},
  backgroundcolor=\color{code-bg},
  frame=single,
  framerule=0.4pt,
  rulecolor=\color{code-border},
  breaklines=true,
  breakatwhitespace=true,
  showstringspaces=false,
  tabsize=4,
  xleftmargin=0.8em,
  xrightmargin=0.8em,
  aboveskip=1em,
  belowskip=1em,
  numbers=none,
  numberstyle=\tiny\color{oreilly-lightgray},
  numbersep=8pt,
  keepspaces=true,
  keywordstyle=\color{oreilly-teal}\bfseries,
  stringstyle=\color{tip-green},
  commentstyle=\color{oreilly-lightgray}\itshape,
  columns=fullflexible,
  literate={€}{{\euro}}1 {→}{{$\rightarrow$}}1 {←}{{$\leftarrow$}}1
}
\lstdefinelanguage{yaml}{
  keywords={true,false,null,yes,no},
  sensitive=false,
  comment=[l]{\#},
  morestring=[b]",
  morestring=[b]',
}

% === INTESTAZIONI E PIÈ DI PAGINA (stile O'Reilly: minimale) ===
\usepackage{fancyhdr}
\pagestyle{fancy}
\fancyhf{}
\fancyhead[LE]{\small\itshape\color{oreilly-gray}\leftmark}
\fancyhead[RO]{\small\itshape\color{oreilly-gray}\rightmark}
\fancyfoot[LE,RO]{\small\thepage}
\renewcommand{\headrulewidth}{0.3pt}
\renewcommand{\headrule}{\hbox to\headwidth{\color{oreilly-lightgray}\leaders\hrule height \headrulewidth\hfill}}

% === TITOLI CAPITOLO (stile O'Reilly: elegante, linea sotto) ===
\usepackage{titlesec}
\titleformat{\chapter}[display]
  {\normalfont\bfseries}
  {\Large\sffamily\color{oreilly-lightgray}Chapter\ \thechapter}
  {8pt}
  {\huge\color{oreilly-dark}}
  [\vspace{4pt}{\color{oreilly-teal}\titlerule[1.5pt]}]
\titlespacing*{\chapter}{0pt}{0pt}{30pt}
\titleformat{\section}
  {\normalfont\Large\bfseries\color{oreilly-dark}}
  {\thesection}{0.8em}{}
\titleformat{\subsection}
  {\normalfont\large\bfseries\color{oreilly-dark}}
  {\thesubsection}{0.8em}{}
\titleformat{\subsubsection}
  {\normalfont\normalsize\bfseries\itshape\color{oreilly-gray}}
  {\thesubsubsection}{0.8em}{}

% === TITOLI PARTI ===
\titleformat{\part}[display]
  {\centering\normalfont\bfseries}
  {\LARGE\color{oreilly-lightgray}\partname\ \thepart}
  {12pt}
  {\Huge\color{oreilly-dark}}
  [\vspace{6pt}{\color{oreilly-teal}\titlerule[2pt]}]

% === HYPERLINK ===
\usepackage{xurl}
\usepackage[
  colorlinks=true,
  linkcolor=oreilly-dark,
  citecolor=oreilly-teal,
  urlcolor=link-blue,
  bookmarks=true,
  bookmarksnumbered=true
]{hyperref}

% === PROTEZIONE OVERFLOW ===
\tolerance=1000
\emergencystretch=1.5em
\setlength{\tabcolsep}{4pt}

% === SPAZIATURA LISTE ===
\usepackage{enumitem}
\setlist{nosep, leftmargin=2em, topsep=0.4em, itemsep=0.2em}

% === BOX ADMONITION ===
% ============================================================
% oreilly-boxes.tex — Box admonition stile O'Reilly
% ============================================================
% Stile: icona a sinistra, bordo sottile, sfondo leggero.
% O'Reilly usa Tip, Note, Warning, Caution, Important.
% ============================================================

\providecolor{tip-green}{HTML}{528B2E}
\providecolor{warn-orange}{HTML}{CC6600}
\providecolor{note-blue}{HTML}{336699}
\providecolor{caution-red}{HTML}{993333}
\providecolor{oreilly-gray}{HTML}{555555}

\usepackage{fontawesome5}
\usepackage{tcolorbox}
\tcbuselibrary{skins,breakable}

% Tip — scorpion icon (O'Reilly classic)
\newtcolorbox{tipbox}[1][]{
  colback=tip-green!3,
  colframe=tip-green!50,
  fonttitle=\sffamily\bfseries\small,
  title={\faLightbulb\space \LabelTip},
  boxrule=0.5pt, arc=2pt, breakable,
  left=8pt, right=8pt,
  coltitle=tip-green!80!black, #1
}

% Note
\newtcolorbox{notebox}[1][]{
  colback=note-blue!3,
  colframe=note-blue!50,
  fonttitle=\sffamily\bfseries\small,
  title={\faStickyNote\space \LabelNote},
  boxrule=0.5pt, arc=2pt, breakable,
  left=8pt, right=8pt,
  coltitle=note-blue!80!black, #1
}

% Warning
\newtcolorbox{warnbox}[1][]{
  colback=warn-orange!4,
  colframe=warn-orange!50,
  fonttitle=\sffamily\bfseries\small,
  title={\faExclamationTriangle\space \LabelWarning},
  boxrule=0.5pt, arc=2pt, breakable,
  left=8pt, right=8pt,
  coltitle=warn-orange!80!black, #1
}

% Caution
\newtcolorbox{dangerbox}[1][]{
  colback=caution-red!3,
  colframe=caution-red!50,
  fonttitle=\sffamily\bfseries\small,
  title={\faTimesCircle\space \LabelCaution},
  boxrule=0.5pt, arc=2pt, breakable,
  left=8pt, right=8pt,
  coltitle=caution-red!80!black, #1
}

% Important
\newtcolorbox{checkpointbox}[1][]{
  colback=tip-green!3,
  colframe=tip-green!40,
  fonttitle=\sffamily\bfseries\small,
  title={\faCheckCircle\space \LabelImportant},
  boxrule=0.5pt, arc=2pt, breakable,
  left=8pt, right=8pt,
  coltitle=tip-green!80!black, #1
}

% Version Note
\newtcolorbox{versionbox}[1][]{
  colback=purple!2,
  colframe=purple!40,
  fonttitle=\sffamily\bfseries\small,
  title={\faCogs\space \LabelVersionNote},
  boxrule=0.5pt, arc=2pt, breakable,
  left=8pt, right=8pt,
  coltitle=purple!70!black, #1
}

% Cost
\newtcolorbox{costbox}[1][]{
  colback=Goldenrod!3,
  colframe=Goldenrod!50,
  fonttitle=\sffamily\bfseries\small,
  title={\faCoins\space \LabelCost},
  boxrule=0.5pt, arc=2pt, breakable,
  left=8pt, right=8pt,
  coltitle=Goldenrod!80!black, #1
}

% Practice
\newtcolorbox{praticabox}[1][]{
  colback=oreilly-gray!3,
  colframe=oreilly-gray!30,
  fonttitle=\sffamily\bfseries\small,
  title={\faWrench\space \LabelExercise},
  boxrule=0.5pt, arc=2pt, breakable,
  left=8pt, right=8pt,
  coltitle=oreilly-gray, #1
}


% === EURO ===
\newcommand{\euro}{€}

% === INDICE ANALITICO ===
\usepackage{makeidx}
\makeindex
